\section{Option Return Model}

\subsection{Local option P\&L}

Let option $i$ be written on underlying $u(i)$ with mark $C_i$, spot
$S_u$, implied volatility state $\sigma_i$, and time-to-expiry $\tau_i$.
For one short allocation horizon, assume the pricing function is twice
differentiable in spot and once differentiable in volatility and time.
Ito--Taylor expansion gives
\begin{equation}
  \Delta C_i
  =
  \Delta_i \Delta S_u
  + \frac12 \Gamma_i(\Delta S_u)^2
  + \nu_i \Delta\sigma_i
  + \Theta_i\Delta t
  + \eta_i,
  \label{eq:taylor_pnl}
\end{equation}
where $\Delta_i$, $\Gamma_i$, $\nu_i$, and $\Theta_i$ are the option's
delta, gamma, vega, and theta at the rebalance date, and $\eta_i$ is the
remainder plus marking error.

Equation \eqref{eq:taylor_pnl} is a portfolio-construction use of the
same second-order term that creates the BSM PDE
\citep{blackscholes1973,merton1973,shreve2004stochasticii}. We do not
assume BSM is true globally. We only use its local Greek accounting.

\subsection{NAV-normalized option returns}

Before introducing Greeks, fix the self-financing accounting.  Let \(W_t\) be beginning
portfolio NAV, let \(q_{i,t}\) be the signed premium weight assigned to option \(i\), and
let \(C_{i,t}>0\) be the option mark.  The position is
\[
  N_{i,t}=\frac{q_{i,t}W_t}{C_{i,t}} .
\]
Ignoring transaction costs and slippage, but not ignoring the option's purchase price, the
mark-to-market return contribution is
\[
  \frac{\Delta W_{i,t+1}}{W_t}
  =
  q_{i,t}\frac{\Delta C_{i,t+1}}{C_{i,t}} .
  \label{eq:premium_accounting}
\]
Equation \eqref{eq:premium_accounting} is the accounting bridge between listed option
marks and portfolio growth.  It rules out free-option P\&L: a long option position whose
mark falls to zero produces \(r_i=-1\) and therefore loses exactly the paid premium
weight \(q_i\).

Define the option return $r_i=\Delta C_i/C_i$.  Write the underlying
simple return as $R_u=\Delta S_u/S_u$.  Dividing
\eqref{eq:taylor_pnl} by $C_i$ gives
\begin{equation}
 r_i
 =
 \underbrace{\frac{\Delta_i S_u}{C_i}}_{b^{\Delta}_{i,u}}R_u
 +\underbrace{\frac12\frac{\Gamma_i S_u^2}{C_i}}_{b^{\Gamma}_{i,u}}R_u^2
 +\underbrace{\frac{\nu_i}{C_i}}_{b^{\nu}_{i,u}}\Delta\sigma_i
 +\underbrace{\frac{\Theta_i\Delta t}{C_i}}_{\theta_i^\ast}
 +\epsilon_i .
 \label{eq:option_return}
\end{equation}
The coefficients are dimensionless option-return exposures.  A deep OTM
option can have small premium $C_i$ and therefore very large return
exposure; this is not a bug but the central reason option-only
Markowitz requires explicit NAV and exposure budgets.

\begin{definition}[Option shock vector]
For $n$ underlyings define
\[
 f=
 (R_1,\ldots,R_n,\;R_1^2-\E R_1^2,\ldots,R_n^2-\E R_n^2,\;
 \Delta\sigma_1,\ldots,\Delta\sigma_n)^\top .
\]
Let $\Omega=\Var(f)$.  For $m$ options, let $B\in\R^{m\times 3n}$ be the
matrix containing the delta-return, gamma-return, and vega-return
loadings in \eqref{eq:option_return}.  Let
$\mu_O=\E r_O$ include theta carry and all expected premia.
\end{definition}

\begin{assumption}[Residual risk]
The residual vector $\varepsilon$ satisfies $\E[\varepsilon]=0$,
$\Cov(f,\varepsilon)=0$, and
$\Sigma_\varepsilon=\Var(\varepsilon)\succeq0$.  In production
estimation, $\Sigma_\varepsilon$ can be diagonal, sparse, shrinkage
cleaned, or estimated from residual option-bucket returns.
\end{assumption}

\begin{proposition}[Greek-induced option covariance]
\label{prop:covariance}
Under the residual-risk assumption,
\[
      \Sigma_O \equiv \Var(r_O) = B\Omega B^\top + \Sigma_\varepsilon .
\]
Moreover, $\Sigma_O$ is positive semidefinite.
\end{proposition}

\begin{proof}
Since $r_O=\mu_O+Bf+\varepsilon$ and constants have zero covariance,
\[
\Var(r_O)=B\Var(f)B^\top+B\Cov(f,\varepsilon)+
\Cov(\varepsilon,f)B^\top+\Var(\varepsilon).
\]
The cross terms vanish by assumption, giving the formula.  For any
$x\in\R^m$,
$x^\top B\Omega B^\top x=(B^\top x)^\top\Omega(B^\top x)\ge0$ because
$\Omega\succeq0$, and $x^\top\Sigma_\varepsilon x\ge0$.  The sum is PSD.
\end{proof}


\subsection{Conditional option-risk-premium means}

The covariance construction says how option returns move together. It does not by itself
create Sharpe. A trader-grade option allocator must identify conditional expected returns
before applying Markowitz. For each option we write
\begin{equation}
 r_{i,t+1}=B_{i,t}f_{t+1}+\theta_{i,t}+\varepsilon_{i,t+1},
\end{equation}
and forecast the conditional mean as
\begin{equation}
\widehat\mu_{i,t}
=
 b^S_{i,t}\widehat\lambda^S_t
+b^\Gamma_{i,t}\widehat\lambda^{var}_t
+b^\nu_{i,t}\widehat\lambda^{vol}_t
+b^{skew}_{i,t}\widehat\lambda^{skew}_t
+b^{tail}_{i,t}\widehat\lambda^{tail}_t
+\theta^\ast_{i,t}.
\label{eq:conditional_mu}
\end{equation}
The terms have economic interpretations. Delta loads on the equity risk premium; gamma
and vega load on variance and volatility risk premia; put skew and VIX call wings load on
crash-insurance and tail-hedging premia; theta is carry under unchanged state variables.
The option-return and variance-risk-premium literatures document why those components can
carry premia rather than merely hedge noise \citep{coval_shumway2001,bakshi_kapadia2003,
carr_wu2009,bollerslev_tauchen_zhou2009,broadie_chernov_johannes2009}. In implementation,
each component is estimated only from training-window or rebalance-date information and
then shrunk toward zero. This shrinkage is part of the model, not an afterthought:
plug-in option means are noisy, and unconstrained mean--variance optimizers otherwise
maximize estimation error.

\begin{proposition}[Conditional option Markowitz]
Given any $\mathcal F_t$-measurable forecasts $\widehat\mu_t$, $B_t$, $\Omega_t$, and
$\Sigma_{\varepsilon,t}\succeq0$, the conditional maximum-Sharpe option problem is
\[
 \max_{q\in\mathcal K_t}
 \frac{q^\top\widehat\mu_t}{\sqrt{q^\top\widehat\Sigma_tq}},
 \qquad
 \widehat\Sigma_t=B_t\Omega_tB_t^\top+\Sigma_{\varepsilon,t}.
\]
If $\mathcal K_t=\mathbb R^m\setminus\{0\}$ and $\widehat\Sigma_t\succ0$, the optimizer is
proportional to $\widehat\Sigma_t^{-1}\widehat\mu_t$.
\end{proposition}

\begin{proof}
Condition on $\mathcal F_t$. The covariance proof in Proposition~\ref{prop:covariance}
applies to the conditional covariance because the forecasts and Greeks are known at $t$.
The tangency proof is identical to Theorem~\ref{thm:tangency} after replacing population
inputs by conditional forecasts.
\end{proof}

\subsection{VIX options as forward-volatility options}

VIX options are not options on an equity spot. The traded forward-risk anchor is the VX
futures curve, while VIX, VVIX, and VIX term-structure indices are state variables. For a
VIX option $j$ expiring at $T_j$, the local return map is therefore
\begin{equation}
 dC^{VIX}_j
 \approx
 \Delta^{VX}_j dF^{VX}_{T_j}
 +\frac12\Gamma^{VX}_j(dF^{VX}_{T_j})^2
 +\nu^{VIX}_j d\sigma^{VIX}_j
 +\Theta^{VIX}_jdt.
\label{eq:vix_forward_map}
\end{equation}
The empirical code computes these Greeks with the Black futures-option model
\citep{black1976commodity}. Expiry P\&L is exact only when the VIX final settlement series
is available for the relevant expiry. The generated settlement audit therefore determines
whether VIX rows are promoted to exact listed-settlement evidence or remain settlement
proxies.

\subsection{Nested models}

The framework nests three useful approximations.
\begin{enumerate}[leftmargin=2em]
  \item \textbf{Delta-only:}
  $\Sigma_O=B_\Delta\Sigma_S B_\Delta^\top+\Sigma_\varepsilon$.
  This is the direct analogue of an equity portfolio whose holdings are
  replaced by option delta-equivalent exposures.
  \item \textbf{Delta-gamma:}
  include the covariance of centered squared returns.  Under joint
  normality,
  $\Cov(R_i^2,R_j^2)=2\Cov(R_i,R_j)^2$.  This makes short-gamma and
  long-gamma books risk differently even when their deltas offset.
  \item \textbf{Delta-gamma-vega:}
  include implied-volatility shocks.  This is the baseline specification
  because equity option books often lose or gain more from vol-surface
  repricing than from spot moves alone.
\end{enumerate}

The portfolio-risk literature emphasizes that production risk models are
factor systems with exposures, factor covariance, and idiosyncratic
covariance, and that optimizers can amplify precision-matrix errors
\citep{ledoitwolf2004}. The option-only model is exactly that risk
model, with Greeks as exposures.
